% Compile with: latexmk -xelatex -interaction=nonstopmode -outdir=. research_proposal_CN.tex
\documentclass[conference]{IEEEtran}
\usepackage{ctex}  % Enables Chinese characters
\usepackage{amsmath, amssymb} % 用于数学公式
\usepackage{graphicx} % 用于插入图像
\usepackage{cite} % 用于管理参考文献
\usepackage{array} % 控制 p{} 列宽
\usepackage{xcolor}
\usepackage{tcolorbox}


\date{\today}

\begin{document}

\title{面向对抗性合作的密码学协议实用框架研究}
\author{
    \IEEEauthorblockN{S. Restrepo Ruiz}
    \IEEEauthorblockA{waajacu.com \\ 电子邮件: santiago.restrepo.ruiz@gmail.com}
}

\maketitle

\begin{abstract}
本研究探讨了在对抗性环境中利用信息理论设计密码学协议以建立信任的方法。
通过使用安全多方计算，本研究旨在构建一个\textit{开源}软件库来展示研究结果。
\end{abstract}

\section{引言}
如果本研究能获得\textit{中华人民共和国工业和信息化部}的奖学金资助，
我们将会开发密码学框架，以期在结构化互动中为合作沟通提供增强。
同时实现信息安全且可验证的加密协议，为博弈论中的协作和隐私保护的决策提供安全保障。

\section{文献综述}
密码学不仅仅是加密技术，还为数字信任奠定了基础。本研究所使用的工具包括：
\textcolor{blue}{\textit{零知识证明（Zero-knowledge proofs}},
\textcolor{blue}{\textit{同态计算（Homomorphic Computation}},
\textcolor{blue}{\textit{混淆电路（Garbled Circuits}},
\textcolor{blue}{\textit{遗忘传输（Oblivious Transfer}},
\textcolor{blue}{\textit{区块链（Blockchains}},
\textcolor{blue}{\textit{差分隐私（Differential Privacy}}，以及
\textcolor{blue}{\textit{门限密码学（Threshold Cryptography}}；
——将这些安全组件组合在一起，能支持广泛的协作式问题求解，这被认为是实现对抗性合作的关键。

\vspace{5pt}
\section{研究问题}

\begin{tabular}{lp{210pt}}
    {\small \textbf{1)}} & {\small 将密码学应用于合作意味着什么？} \\
    {\small \textbf{2)}} & {\small 如何通过跨学科合作来改进所提出的密码学协议？}
\end{tabular}

\vspace{5pt}
\section{研究目标}

\begin{tabular}{lp{210pt}}
    {\small \textbf{1)}} & {\small 在经典博弈论问题上进行验证演示。} \\
    {\small \textbf{2)}} & {\small 基于此研究成果开源一套软件库。}
\end{tabular}

\section{研究方法}
\vspace{5pt}
本研究将主要关注以下结构化博弈场景：
\vspace{5pt}

\begin{tabular}{lp{210pt}}
    {\small \textbf{1)}} & {\textcolor{blue}{\textit{在不暴露手牌的前提下玩纸牌游戏：}} \linebreak---演示如何在纸牌游戏中宣布获胜，而不公开具体手牌。 } \\
    {\small \textbf{2)}} & {\textcolor{blue}{\textit{在不暴露策略的前提下玩井字棋：}} \linebreak---演示玩家如何在不公开具体落子位置的情况下，证明自己拥有取胜策略。} \\
    {\small \textbf{3)}} & {\textcolor{blue}{\textit{经典对抗性“捕鹿”博弈（stag-hunt）：}} \linebreak---演示在这一对抗性合作的基本博弈中，如何找到通往均衡的路径。}
\end{tabular}

\section{资助申请及理由}
我诚恳地请求\textit{北京航空航天大学}的各位评审专家对我的申请给予支持。 

目前在哥伦比亚的大学中，尚缺乏这一领域的研究项目。我已经坚持学习中文多年，
深知在中国学术环境中顺畅的沟通与协作的重要性。
我对能在中国这一充满活力的学术社区中学习和做出贡献充满热情。
我相信，这项奖学金不仅能帮助我实现职业抱负，也能促进我们两国间宝贵的学术交流。

\section{目标大学及相关项目}
\vspace{5pt}
{\small 英文授课项目:} \\

\begin{tabular}{rl}
    \textcolor{blue}{\small\textit{北京航空航天大学}} & {\small 硕士学位：网络空间安全}\\
    \textcolor{blue}{\small\textit{北京航空航天大学}} & {\small 硕士学位：集成电路工程}
\end{tabular}

\section{预期成果}

\begin{tabular}{lp{210pt}}
    {\small \textbf{1)}} & {\small 在博弈论中演示对抗性合作。} \\
    {\small \textbf{2)}} & {\small 开发一套在对抗性环境中实现合作的开源密码协议套件。} \\
    {\small \textbf{3)}} & {\small 倡导使用密码学在冲突中建立信任的可行性。}
\end{tabular}

\section{结论}
本研究旨在展示密码学技术如何应用于和平构建与信任建立。

% \begin{thebibliography}{99}

% \bibitem{Yao1982}
% A. C. Yao, ``Protocols for secure computations,'' in *Proceedings of the 23rd Annual Symposium on Foundations of Computer Science*, 1982.

% \bibitem{Goldwasser1989}
% S. Goldwasser, S. Micali, and C. Rackoff, ``The Knowledge Complexity of Interactive Proof Systems,'' in *SIAM Journal on Computing*, 1989.

% \bibitem{Gentry2009}
% C. Gentry, ``Fully Homomorphic Encryption Using Ideal Lattices,'' in *Proceedings of the 41st ACM Symposium on Theory of Computing (STOC)*, 2009.

% \bibitem{Dwork2006}
% C. Dwork, ``Differential Privacy,'' in *Proceedings of the International Colloquium on Automata, Languages, and Programming (ICALP)*, 2006.

% \bibitem{Canetti2000}
% R. Canetti, ``Security and Composition of Multi-Party Cryptographic Protocols,'' in *Journal of Cryptology*, 2000.

% \end{thebibliography}

% \bibliographystyle{IEEEtran}
% \bibliography{references}

\end{document}
